This chapter examines how AWS Lambda complements edge computing within a distributed architecture. The focus is not on Lambda as an edge technology, but on how its event-driven execution model aligns with the operational logic of edge systems and supports the cloud-side components of an edge–cloud continuum.

Edge applications typically generate irregular, context-dependent events that require cloud-side handling without maintaining long-running servers. Lambda fits naturally into this pattern because it executes only when triggered \cite{aws_lambda_runtime, aws_serverless_whitepaper}. In our Edge Computing report, anomaly detections were forwarded to the cloud via MQTT, illustrating how edge devices act as event producers \cite{edge_computing_report}. AWS IoT Core and its rule engine transform such events into Lambda invocations, resulting in a loosely coupled architecture in which the edge reacts to sensor data while Lambda reacts to the outputs of the edge. This makes both layers lightweight and autonomous.

Lambda is particularly suitable as a cloud-side control plane for edge deployments. Many responsibilities that exceed the scope of a single edge node, such as aggregating anomalies, enriching events with additional context or initiating workflow logic, are better handled centrally \cite{aws_serverless_whitepaper}. Lambda performs these tasks without persistent servers and scales automatically in response to inputs from many distributed nodes. This model complements the decentralised edge processing described in our Edge Computing report, providing global coordination without interfering with local autonomy \cite{edge_computing_report}.

However, not every edge workload is appropriate for Lambda. A key limitation is cold starts, a delay occurring when AWS must initialise a fresh execution environment because no warm instance is available. The Lambda report \cite{lambda_report} demonstrates this behaviour and analyses mitigation mechanisms such as SnapStart \cite{aws_snapstart} and Provisioned Concurrency \cite{aws_provisioned_concurrency}, both intended to reduce cold-start latency. Academic work confirms that cold starts cause variability when invocations are sporadic or bursty, which is typical of many edge workloads \cite{wang_peeking_serverless}. As a result, real-time inference or time-critical loops should remain at the edge, while Lambda is suited for asynchronous or moderately time-sensitive cloud-side tasks.

Lambda’s support for Response Streaming, also demonstrated in the Lambda report \cite{lambda_report} and specified in AWS documentation \cite{aws_response_streaming}, improves cloud–edge interaction in scenarios requiring timely feedback. Response Streaming allows clients to receive partial outputs before the full computation completes, reducing perceived latency. For edge applications that depend on dashboards, operator updates or progressive anomaly reporting, such incremental feedback can provide valuable situational awareness even before final results are available.

Finally, Lambda should not be seen as a substitute for edge computation. Its stateless execution model \cite{aws_lambda_runtime, aws_serverless_whitepaper}, lack of hardware acceleration \cite{wang_peeking_serverless} and execution-time limits make it unsuitable for continuous or inference-heavy workloads. Instead, Lambda’s strengths lie in elasticity, large-scale event processing and central orchestration. Used together, edge compute provides locality and determinism while Lambda provides scalable serverless coordination. This complementary relationship enables distributed systems that are efficient at the periphery yet manageable and extensible at the cloud layer.
